\documentclass[11pt]{article}
\usepackage[margin=1in]{geometry}
\usepackage{hyperref}
\title{TLA+ Study Group — Week 01}
\author{@benetis}
\date{\today}
\begin{document}\maketitle

\section{Questions}

\subsection{\textit{if} ($\Rightarrow$) and \textit{iff} ($\Leftrightarrow$)}

\subsubsection*{Definition}
\begin{enumerate}
  \item $\Leftrightarrow$ is the logical equivalence.
  \item $p \Leftrightarrow q$ is read as \underline{p equivalent q} or \underline{p if and only if q}
\end{enumerate}

\subsubsection*{Truth table}

\begin{tabular}{c|c|c}
p & q & p $\Leftrightarrow$ q \\
\hline 
T & T & T \\
T & F & F \\
F & T & F \\
F & F & T \\
\end{tabular}

\subsubsection*{Example}
\textit{if} ($\Rightarrow$) is one-way: If a number is divisible by 4, then it is even. (But not all even numbers are divisible by 4.)

\noindent\textit{iff} ($\Leftrightarrow$) must hold both ways: A number is even $\Leftrightarrow$ it is divisible by 2. (Even $\Rightarrow$ divisible by 2, and divisible by 2 $\Rightarrow$ even.)

\noindent\textbf{Note:} An example of this in TLA+ can be found in \texttt{tla\_plus/iff.tla}


\end{document}
